% Paper draft — see draft1/ for the bundled version with figures/ and tables/.
% Compile: cd final/paper/draft1 && pdflatex draft1.tex
\documentclass[11pt]{article}
\usepackage[margin=1in]{geometry}
\usepackage{amsmath,amssymb}
\usepackage{graphicx}
\graphicspath{{./}{figs/}{figures/}}
\usepackage{booktabs}
\usepackage{caption}
\usepackage{subcaption}
\usepackage{natbib}
\usepackage{placeins}
\usepackage[colorlinks=true,allcolors=blue]{hyperref}

\title{Reward-Proportional Sampling in Large DAGs with\\
Multiplicity-Corrected Inverse Probability Scaling}
\author{}
\date{}

\begin{document}
\maketitle

%======================================================================
\section{Introduction}
%======================================================================

Many reinforcement learning problems have more than one good answer. In those
settings we do not want the single best terminal outcome, we want to sample
terminal outcomes in proportion to their reward,
\begin{equation}
p^{*}(x) \;=\; \frac{R(x)}{Z}, \qquad Z = \sum_{x} R(x).
\label{eq:target}
\end{equation}
Policies trained by expected-return maximization do not do this. They collapse
onto a small set of outcomes, and this collapse comes from the objective itself
rather than from weak exploration \citep{sinha2026ips}. Inverse probability
scaling (IPS) removes the frequency multiplier from the learning signal by
dividing the terminal reward by the outcome probability, and IPS-GRPO
implements this by estimating that probability with within-group counts.

Counting works when the outcome space is small enough that the same outcome
appears several times in one group. It stops working in two situations we care
about. First, when the number of reachable outcomes is much larger than the
group size, every sampled outcome is unique and $\hat{p}(x)=1/G$ for all of
them, so the correction becomes a constant and IPS-GRPO reduces to GRPO.
Second, when the environment is a DAG rather than a tree, one outcome is
reachable by many different trajectories, and a trajectory-level weight has to
account for that.

This paper replaces the count estimate with a learned reverse policy
$q_\phi(\tau \mid x)$, and calls the resulting algorithm \emph{multiplicity-corrected
inverse probability scaling} (MIPS-GRPO). The reverse policy gives a
trajectory-level weight that is unbiased at the outcome level for any DAG, and
does not depend on seeing repeats inside a group.

Correcting for multiplicity is not by itself new: a GFlowNet with a fixed
uniform backward policy $P_B$ already does it. The contribution here is that
$q_\phi$ is fitted to $P_F(\cdot \mid x)$ rather than left uniform. Any
normalized reverse policy gives the same terminal distribution, so fitting
cannot change what the method converges to, but it controls the variance of the
weights and therefore how fast it gets there. That distinction is invisible on
small problems and decisive on large ones: at 27 taxa a uniform reverse policy
leaves the effective sample size at $0.847$, while a fitted one reaches $0.999$,
and the difference between the two is the difference between a usable sampler
and an unusable one.

We test the method on a $64\times64$ hypergrid at group size $32$, where
counting is provably uninformative, and then on phylogenetic tree sampling with
up to $27$ taxa, where the outcome space cannot be enumerated at all.

%======================================================================
\section{Setup}
%======================================================================

The environment is a DAG. A forward policy $P_F$ starts at a fixed root and
produces a trajectory $\tau = (s_0, a_1, \dots, a_T, s_T)$ ending at a terminal
outcome $x = s_T$. The forward policy gives the probability of a path,
$P_F(\tau) = \prod_t P_F(a_t \mid s_{t-1})$, but the objective in
Eq.~\eqref{eq:target} is defined over terminals, and
\begin{equation}
P_F(x) \;=\; \sum_{\tau \in \mathcal{T}(x)} P_F(\tau)
\end{equation}
requires summing over every trajectory that ends at $x$. That sum is computable
in a small toy DAG and nowhere else.

\paragraph{The naive weight is wrong.}
Since $P_F(\tau)$ is known exactly for a sampled rollout, it is tempting to use
$w = R(x)/P_F(\tau)$. But then every trajectory ending at $x$ contributes
$P_F(\tau) \cdot R(x)/P_F(\tau) = R(x)$, so a terminal reachable by
$m(x) = |\mathcal{T}(x)|$ paths receives target mass $m(x)R(x)$ instead of
$R(x)$. Terminals with more paths are systematically over-weighted. This is not
a small effect. On the $64\times64$ hypergrid of Section~\ref{sec:grid}, the
four modes carry identical reward but their multiplicities span about
$7\times10^{5}$ to $1.6\times10^{30}$.

\paragraph{The fix.}
Introduce a distribution $q_\phi(\tau \mid x)$ over the trajectories that could
have produced $x$, normalized per terminal,
$\sum_{\tau \in \mathcal{T}(x)} q_\phi(\tau \mid x) = 1$. The corrected weight
is
\begin{equation}
w(\tau) \;=\; \frac{R(x)\, q_\phi(\tau \mid x)}{P_F(\tau)},
\label{eq:weight}
\end{equation}
and now the paths to $x$ sum to the right thing:
\begin{equation}
\sum_{\tau \in \mathcal{T}(x)} P_F(\tau)\, w(\tau)
= R(x) \sum_{\tau \in \mathcal{T}(x)} q_\phi(\tau \mid x) = R(x).
\end{equation}
The choice of $q_\phi$ moves mass between the paths to a terminal but never
changes the total mass of that terminal.

\paragraph{Why learn $q_\phi$ instead of fixing it uniform.}
Any normalized $q_\phi$ is unbiased, so a fixed uniform reverse policy (what
GFlowNet uses as $P_B$) is already valid. The issue is variance. If
$q_\phi = P_F(\tau \mid x)$ exactly, then
\begin{equation}
w(\tau) = \frac{R(x) P_F(\tau \mid x)}{P_F(x) P_F(\tau \mid x)} = \frac{R(x)}{P_F(x)},
\end{equation}
which is constant across all paths to $x$: zero within-terminal variance, and
the weight equals the ideal terminal-level weight we could not compute
directly. The further $q_\phi$ is from $P_F(\cdot \mid x)$, the more the weights
scatter and the lower the effective sample size (ESS). Normalization is what
makes $q_\phi$ correct. Learning $q_\phi$ is what makes it low variance. A badly
fit $q_\phi$ hurts ESS but cannot bias the terminal distribution.

%======================================================================
\section{Method}
%======================================================================

Each edge of a sampled trajectory is one labelled example: given the child state
and the terminal, predict which incoming action was taken,
\begin{equation}
\text{input } (s_{i,t}, x_i) \;\longrightarrow\; \text{target } a_{i,t}.
\end{equation}
The reverse network outputs a masked categorical over the valid parent edges of
$s_{i,t}$, and the path probability is the product over steps,
$\log q_\phi(\tau_i \mid x_i) = \sum_t \log q_\phi(a_{i,t} \mid s_{i,t}, x_i)$.
Training is maximum likelihood on the on-policy batch:
\begin{equation}
\mathcal{L}_{\text{rev}}(\phi) = -\frac{1}{N} \sum_{i=1}^{N} \sum_{t=1}^{T_i}
\log q_\phi(a_{i,t} \mid s_{i,t}, x_i).
\label{eq:revloss}
\end{equation}
Rewards do not appear here. The reverse model only learns which paths the
forward policy tends to use given the terminal, which is exactly
$P_F(\tau \mid x)$ in the realizable limit.

\paragraph{One iteration.}
\begin{enumerate}\itemsep2pt
\item Sample trajectories with $P_F$; record log-probs, actions, terminals, rewards.
\item Score them with the \emph{frozen} reverse model to get $q_\phi(\tau_i \mid x_i)$.
\item Form $\log w = \log R(x) + \log q_\phi(\tau \mid x) - \log P_F(\tau)$,
normalize to advantages with a running EMA, and update $P_F$ with path-level PPO.
\item Fit $q_\phi$ on the same batch by minimizing Eq.~\eqref{eq:revloss}; use the
updated model for the next batch.
\end{enumerate}
The freeze-then-fit order matters. Fitting before scoring would inflate
$q_\phi$ on whichever paths happened to appear, which is self-reinforcing. The
reverse model never generates rollouts, it only scores them.

\paragraph{Comparison to the baselines.}
\begin{center}
\begin{tabular}{lll}
\toprule
Method & Training signal & Multiplicity? \\
\midrule
GRPO & group-normalized rewards $\to$ PPO & no \\
IPS-GRPO & $R(x)/\hat{p}(x)$, $\hat{p}$ from batch counts $\to$ PPO & no \\
GFlowNet TB & squared loss, fixed uniform $P_B$, learned $\log Z$ & yes \\
MIPS-GRPO & $R(x) q_\phi(\tau\mid x)/P_F(\tau) \to$ PPO & yes \\
\bottomrule
\end{tabular}
\end{center}
At initialization our $q_\phi$ is uniform, that is, identical to GFlowNet's
$P_B$. The difference is that it is then fit to $P_F(\cdot \mid x)$, and that
$Z$ is absorbed by the advantage normalization instead of being learned as a
separate scalar.

%======================================================================
\section{Hyper-Grid Task}
\label{sec:grid}
%======================================================================

\paragraph{Task setup.}
We use the two-dimensional hyper-grid of \citet{sinha2026ips}. Episodes start at
$x_0 = (0,0)$, and an action either increments one coordinate by $1$ or
terminates the episode and collects the terminal reward. Terminals are grid
coordinates $x \in \{0,\dots,H-1\}^2$, and the reward is
\begin{equation}
\begin{aligned}
R(x) ={}& R_0 + R_1 \prod_{i=1}^{n}\mathbf{1}\bigl(|x_i| > 0.5\bigr) \\
        &+ R_2 \prod_{i=1}^{n}\mathbf{1}\bigl(0.6 < |x_i| < 0.8\bigr),
\end{aligned}
\label{eq:gridreward}
\end{equation}
where $n=2$ in our experiments, $0 < R_0 \ll R_1 < R_2$, and coordinates are
normalized to $[0,1]$. This gives
four symmetric high-reward modes inside a lower ring
(Figure~\ref{fig:grid-heatmaps}a). We use $R_0 = 0.1$, $R_1 = 0.5$,
$R_2 = 2.0$, so the peak reward is $2.6$. The target distribution is
$p(x) \propto R(x)$ and it is known in closed form, so we can measure the
$\ell_1$ distance between the empirical terminal distribution and the target
exactly.

\paragraph{Why we run it at $H=64$ and $G=32$.}
\citet{sinha2026ips} report this task at a grid size where counting works. We
run the same reward landscape at $H = 64$, which gives $4096$ terminal states,
with group size $G = 32$. Two things break in this regime, and they are the two
things the reverse policy is meant to fix.

First, the count estimate is uninformative. With $4096$ terminals and $32$
samples per group, the sampled outcomes in a group are almost always distinct,
so $\hat{p}(x) = 1/G$ for every one of them. The scaling $R(x)/\hat{p}(x)$ is
then a constant multiple of $R(x)$, and IPS-GRPO becomes GRPO with a rescaled
reward. Clipping at $\epsilon$ does not help, because the problem is not that
$\hat{p}$ is small but that it is the same for everything.

Second, path multiplicity is extreme. A terminal $(a,b)$ is reachable by
$\binom{a+b}{a}$ distinct trajectories. The four modes sit near $(11,11)$,
$(11,52)$, $(52,11)$ and $(52,52)$, with multiplicities of about
$7.1\times10^{5}$, $6.2\times10^{11}$, $6.2\times10^{11}$ and
$1.6\times10^{30}$. The modes carry identical reward and differ in reachability
by roughly twenty-four orders of magnitude. Any trajectory-level weight that
ignores multiplicity will put almost all of its mass on the far corner.

\paragraph{Baselines and metrics.}
We compare against GRPO, count-based IPS-GRPO, and a GFlowNet trained with the
trajectory balance objective under a fixed uniform $P_B$. All four methods share
the architecture, optimizer, sampling budget and entropy regularization. We
report four numbers: $\ell_1$ distance to the target, the number of the four
modes that are visited, the number of distinct terminals seen out of $4096$, and
mode recovery as a function of samples drawn from the trained policy.

\paragraph{Results.}
Table~\ref{tab:grid} and Figure~\ref{fig:grid-heatmaps} give the summary. GRPO
and count IPS-GRPO both collapse onto a single mode and end at $\ell_1$ of
$1.966$ and $1.844$. They also visit almost nothing: $15$ and $43$ distinct
terminals out of $4096$. The small gap between them is what the argument above
predicts. Once every outcome in a group is unique, the count correction has
nothing left to correct, and count IPS-GRPO inherits GRPO's behaviour.

MIPS-GRPO and GFlowNet TB both recover all four modes, reach $\ell_1$
of $0.325$ and $0.323$, and cover $4023$ and $4067$ terminals. These two are
tied on this task, and we do not claim otherwise. The point of the hyper-grid
here is not to separate MIPS-GRPO from GFlowNet TB, it is to show that
the learned reverse weight recovers the reward-proportional distribution in a
setting where the count-based estimate provably cannot, and to check it against
a method whose correctness is already established. The separation between the
two appears in Section~\ref{sec:phylo}, where the outcome space is large enough
that a uniform $P_B$ becomes a poor proposal.

Figure~\ref{fig:modes} shows mode discovery against samples drawn. Both
methods that correct for multiplicity find all four modes within roughly $10^2$ samples,
while GRPO and count IPS-GRPO stay at one mode for the full $5\times10^4$
sample budget. Extra sampling does not rescue them, which is the signature of
outcome-level collapse rather than of slow exploration.

Figure~\ref{fig:train-reward} is worth reading alongside the table. GRPO and
count IPS-GRPO reach mean terminal reward $2.6$, the maximum available, within a
few thousand steps and stay there. MIPS-GRPO and GFlowNet TB settle
near $1.0$. The two collapsed methods are winning on the metric they optimize.
Mean reward is the wrong metric for this problem, and a method that samples in
proportion to reward should be expected to score worse on it.

\begin{table}[htbp]
\centering
\caption{Hyper-grid, $H=64$, group size $32$, four modes. Recovery is the
fraction of modes visited. $\ell_1$ is the distance between the empirical
terminal distribution and $p(x) \propto R(x)$; lower is better. Unique terminals
is out of $4096$.}
\label{tab:grid}
\begin{tabular}{lrrrr}
\toprule
Method & Recovery (\%) & Modes & $\ell_1$ & Unique terminals \\
\midrule
GRPO                     & 25.0  & 1 & 1.966 & 15 \\
Count IPS-GRPO           & 25.0  & 1 & 1.844 & 43 \\
MIPS-GRPO      & 100.0 & 4 & 0.325 & 4023 \\
GFlowNet TB              & 100.0 & 4 & 0.323 & 4067 \\
\bottomrule
\end{tabular}
\end{table}

\begin{figure}[htbp]
\centering
\includegraphics[width=\textwidth]{terminal_distribution_comparison.png}
\caption{Terminal distributions on the $64\times64$ hyper-grid. (a) target
$p(x) \propto R(x)$. (b) GRPO and (c) count IPS-GRPO put nearly all mass on one
of the four equal-reward modes. (d) MIPS-GRPO and (e) GFlowNet TB
cover all four and reproduce the ring.}
\label{fig:grid-heatmaps}
\end{figure}

\begin{figure}[htbp]
\centering
\begin{subfigure}{0.49\textwidth}
  \includegraphics[width=\textwidth]{modes_found_vs_samples_logx.png}
  \caption{log scale}
\end{subfigure}
\hfill
\begin{subfigure}{0.49\textwidth}
  \includegraphics[width=\textwidth]{modes_found_vs_samples.png}
  \caption{linear scale}
\end{subfigure}
\caption{Modes found against samples drawn from the trained policy. The two
methods that correct for multiplicity reach all four modes within about $10^2$ samples.
GRPO and count IPS-GRPO stay at one mode for the whole budget.}
\label{fig:modes}
\end{figure}

\begin{figure}[htbp]
\centering
\begin{subfigure}{0.49\textwidth}
  \includegraphics[width=\textwidth]{training_l1_curve.png}
  \caption{evaluation $\ell_1$}
\end{subfigure}
\hfill
\begin{subfigure}{0.49\textwidth}
  \includegraphics[width=\textwidth]{training_reward_curve.png}
  \caption{training reward}
\end{subfigure}
\caption{Left: $\ell_1$ to the target during training. The two collapsed methods
never improve. Right: mean terminal reward per batch. GRPO and count IPS-GRPO
saturate at the maximum reward $2.6$ while sampling one mode; the
reward-proportional methods sit near $1.0$ by design.}
\label{fig:train-reward}
\end{figure}

\FloatBarrier

%======================================================================
\section{Application: phylogenetic sampling}
\label{sec:phylo}
%======================================================================

A tree is built by repeatedly merging subtrees. At each of the $N-1$ steps the
forward policy picks a pair to merge, from $\binom{k}{2}$ options, and then
branch lengths for the new edge. The reward is a shifted log-likelihood
$R(x) = c + \log L(x)$, with $c$ chosen so that $R > 0$, which the importance
weights require. Evaluation is at \emph{signature} level: an outcome is a
(topology, branch lengths) pair, which is finer than topology alone and matches
the full generative model. All methods draw 1M samples from the trained policy,
and the primary metric is the Pearson correlation between estimated terminal
probability and reward.

\begin{center}
\begin{tabular}{lccc}
\toprule
 & 5 taxa & 10 taxa & 27 taxa \\
\midrule
Dataset        & DS1 reduced & DS1 reduced & DS1 (full) \\
Reward shift $c$ & 3600 & 5000 & 12000 \\
Topologies     & 105 & $\sim3.4\times10^{7}$ & very large \\
Batch / epochs & 4096 / 10k & 4096 / 10k & 1024 / 32k \\
Replay         & none & none & none \\
\bottomrule
\end{tabular}
\end{center}

\begin{table}[h]
\centering
\caption{Phylogenetic sampling results, 1M samples per method.}
\begin{tabular}{llrrr}
\toprule
Setting & Method & Pearson $r$ & Unique sig. / 1M & Coverage \\
\midrule
5 taxa & \textbf{MIPS-GRPO} & \textbf{0.994} & 951{,}175 & 105/105 topologies \\
       & GFlowNet & 0.982 & 960{,}850 & 105/105 \\
       & IPS-GRPO & 0.319 & 66{,}268 & 3/105 \\
       & GRPO & collapsed & 1 & 1/105 \\
\midrule
10 taxa & \textbf{MIPS-GRPO} & \textbf{0.976} & 999{,}986 & $R^2 = 0.952$ \\
        & GFlowNet & 0.881 & 1{,}000{,}000 & $R^2 = 0.777$ \\
        & IPS-GRPO & 0.050 & 140{,}109 & $R^2 = 0.003$ \\
        & GRPO & 0.033 & 2{,}850 & $R^2 = 0.001$ \\
\midrule
27 taxa & \textbf{MIPS-GRPO} & \textbf{0.977} & 999{,}987 & ESS 0.999 \\
        & GFlowNet & 0.002 & 1{,}000{,}000 & ESS 0.847 \\
\bottomrule
\end{tabular}
\end{table}

At 5 taxa MIPS-GRPO and GFlowNet both recover the ideal distribution and
cover all 105 topologies, while IPS-GRPO finds 3 topologies and GRPO finds 1. At
10 taxa the learned reverse policy is ahead on correlation at equal coverage,
and it gets there faster: at epoch 1000 it is already at $r = 0.953$ against
$0.861$ for GFlowNet. At 27 taxa the gap is the largest result in the paper:
near-perfect correlation (topology TV $0.010$, log-weight std $0.027$) against
no meaningful probability-reward relationship for GFlowNet, whose TB loss
($0.430$) and $\log Z$ ($279.65$) were nonetheless stable and decreasing.

Figures~\ref{fig:phylo-5} and~\ref{fig:phylo-10} show the probability--reward
relationship for all four methods. At 5 taxa, both multiplicity-aware methods
track the target, whereas the count-based methods cover only a small fraction
of the outcome space. At 10 taxa, MIPS-GRPO remains tightly aligned with the
ideal linear relationship while the other methods show substantially more
scatter or collapse.

\begin{figure}[htbp]
\centering
\includegraphics[width=0.96\textwidth]{5taxa/sampling_comparison.png}
\caption{Five-taxa outcome probability versus terminal reward from 1M samples.
GRPO collapses and count IPS-GRPO covers few signatures; GFlowNet TB and
MIPS-GRPO both recover the reward-proportional relationship, with MIPS-GRPO
giving the strongest fit.}
\label{fig:phylo-5}
\end{figure}

\begin{figure}[htbp]
\centering
\includegraphics[width=0.96\textwidth]{10taxa/sampling_comparison.png}
\caption{Ten-taxa outcome probability versus terminal reward from 1M samples.
MIPS-GRPO closely follows the ideal relationship, while GFlowNet TB has greater
dispersion and GRPO and count IPS-GRPO have almost no probability--reward
correlation.}
\label{fig:phylo-10}
\end{figure}

Figure~\ref{fig:phylo-27} isolates the largest-scale comparison. Despite both
methods producing almost entirely unique signatures, MIPS-GRPO preserves a
strong probability--reward relationship and GFlowNet TB does not.

\begin{figure}[htbp]
\centering
\begin{subfigure}[t]{0.49\textwidth}
  \includegraphics[width=\textwidth]{27taxa/sampling.png}
  \caption{MIPS-GRPO}
\end{subfigure}
\hfill
\begin{subfigure}[t]{0.49\textwidth}
  \includegraphics[width=\textwidth]{27taxa/gflownet_sampling.png}
  \caption{GFlowNet TB}
\end{subfigure}
\caption{Twenty-seven taxa on full DS1, evaluated with 1M samples. MIPS-GRPO
achieves Pearson $r=0.977$, whereas GFlowNet TB has no meaningful
probability--reward relationship ($r=0.002$).}
\label{fig:phylo-27}
\end{figure}

\paragraph{The reverse policy is the bottleneck at scale.}
At 27 taxa with 10k epochs the reverse loss was still $4.82$ and ESS only
$0.49$, and sampling quality was correspondingly poor. Extending to 32k epochs
with 8 reverse MLE steps per update drove the reverse loss to
$\sim4\times10^{-5}$ and ESS to $0.999$. This is the clearest evidence that the
fit of $q_\phi$, and not the forward objective, controls how well the method
works. Figure~\ref{fig:phylo-training} shows the corresponding optimization
curves across all three problem sizes.

\begin{figure}[htbp]
\centering
\includegraphics[width=\textwidth]{training_all_taxa.png}
\caption{MIPS-GRPO training diagnostics for 5, 10, and 27 taxa. Columns show
reverse-policy negative log-likelihood, IPS effective sample size, and mean log
score. The dashed line marks 10k epochs for the 27-taxa run; continued reverse
training raises ESS toward one and stabilizes the sampling objective.}
\label{fig:phylo-training}
\end{figure}

\begin{center}
\begin{tabular}{lrrrr}
\toprule
Taxa & Epochs & Final reverse loss & Final IPS ESS & Final mean log score \\
\midrule
5  & 10{,}000 & $2.04\times10^{-5}$ & 0.9998 & 267.0 \\
10 & 10{,}000 & $6.46\times10^{-5}$ & 0.9944 & 335.6 \\
27 & 10{,}000 & $4.82$              & 0.488  & 587.3 \\
27 & 32{,}000 & $3.87\times10^{-5}$ & 0.999  & 2026.2 \\
\bottomrule
\end{tabular}
\end{center}

\FloatBarrier

\bibliographystyle{plainnat}
\begin{thebibliography}{9}
\bibitem[Sinha et al.(2026)]{sinha2026ips}
A. Sinha, S. Elango, and D. Liu.
\newblock Expected return causes outcome-level mode collapse in reinforcement
learning and how to fix it with inverse probability scaling.
\newblock arXiv:2601.21669, 2026.
\end{thebibliography}

\end{document}